\subsubsection{Why Distributed Systems Are Special}\label{subsubsec:introduction--why-distributed-systems--why-distributed-systems-are-special}

\begin{flushleft}
    \textcolor{Green3}{\faIcon{balance-scale} \textbf{Distributed vs. Centralized Systems}}
\end{flushleft}
Start from the simplest possible centralized program: an \textbf{application running on a single machine}. Its components share the same machine, operating system, memory hierarchy, and physical environment.If one part of the program invokes another function, conceptually we have:
\begin{equation*}
    A \xrightarrow{\text{function call}} B
\end{equation*}
And both $A$ and $B$ are local (i.e., on the same machine).

\highspace
Now \textbf{distribute} those components $A$ and $B$ with a network between them:
\begin{equation*}
    A \xrightarrow{\text{network}} B
\end{equation*}
That apparently small change is \hl{enormously significant}. A local interaction may involve something like a function call, while a distributed interaction requires:
\begin{enumerate}
    \item Serialization of data into a message
    \item Communication over a network
    \item Network transmission, which may be delayed or lost
    \item Reception of the message by $B$
    \item Deserialization of the message into data
    \item Processing of the data by $B$
\end{enumerate}
And every one of those steps introduces additional uncertainty and potential failure. This is why \hl{a distributed system is not simply a centralized program running on several machines.} The network becomes part of the computation.

\highspace
\begin{flushleft}
    \textcolor{Red2}{\faIcon{question-circle} \textbf{Why building distributed systems is harder}}
\end{flushleft}
As we saw in the previous comparison, distributed systems are more complex than centralized systems. Not only is \textbf{building them harder}, but \textbf{building them \underline{correctly} is much harder}.

\highspace
In general, it is often relatively easy to make a distributed application work when everything behaves normally. For example:
\begin{equation*}
    A \to B, \quad B \to A
\end{equation*}
Request arrives, reply arrives, everything is fine. The \textbf{difficulty} appears when reality interferes. Suppose $A$ sends a request to $B$:
\begin{equation*}
    \texttt{transfer(100 \euro)}
\end{equation*}
$B$ receives the request and performs the operation. Then $B$ sends a reply back to $A$, but it is lost in the network. \emph{What should $A$ do?}
\begin{itemize}
    \item If it sends the request again, perhaps the transfer happens \textbf{twice}.
    \item If it does not retry, perhaps the first request had actually been lost and the transfer happens \textbf{zero times}.
\end{itemize}
So even something that sounds trivial, like ``\emph{retry when there is an error}'', may not be trivial at all. The \hl{problem is that different machines have only \textbf{partial knowledge} of what happened.}

\highspace
\begin{flushleft}
    \textcolor{Green3}{\faIcon{book} \textbf{The Eight Fallacies of Distributed Computing}}
\end{flushleft}
The \definition{Eight Fallacies of Distributed Computing}, associated with Peter Deutsch, are \hl{assumptions that programmers are tempted to make when first designing distributed software}:
\begin{equation*}
    \boxed{
        \begin{array}{ll}
            1.&\text{The network is reliable}\\[.3em]
            2.&\text{Latency is zero}\\[.3em]
            3.&\text{Bandwidth is infinite}\\[.3em]
            4.&\text{The network is secure}\\[.3em]
            5.&\text{Topology doesn't change}\\[.3em]
            6.&\text{There is one administrator}\\[.3em]
            7.&\text{Transport cost is zero}\\[.3em]
            8.&\text{The network is homogeneous}
        \end{array}
    }
\end{equation*}
They are called ``\emph{fallacies}'' because \textbf{none of these assumptions is generally safe}. Let's analyze each of them in turn:
\begin{enumerate}
    \item \important{The network is reliable}. A common mental model is:
        \begin{equation*}
            A \xrightarrow{\text{message}} B
        \end{equation*}
        Therefore $B$ receives the message. But in reality, this is not necessarily true. Packets can be lost, corrupted, delayed, duplicated, or delivered after a connection failure. So the actual model is closer to:
        \begin{equation*}
            A \xrightarrow[\text{maybe}]{\text{network}} B
        \end{equation*}
        This means \hl{distributed software often needs mechanisms such as acknowledgments, retries, duplicate detection, or replication}. However, the important principle is:
        \begin{equation*}
            \text{sending} \neq \text{receiving}
        \end{equation*}

    \item \important{Latency is zero}. A local memory access or function call is extremely fast compared with network communication. If $A$ communicates with $B$:
        \begin{equation*}
            A \xrightarrow[\Delta t]{}
        \end{equation*}
        Where $\Delta t$ is the time it takes for the message to travel from $A$ to $B$, which is $\Delta t > 0$, and it may vary substantially.

        \hl{A remote operation can therefore never safely be treated exactly like an ordinary local operation.} For example, imagine writing:
    \begin{lstlisting}
for each item:
    remoteServer.get(item)
\end{lstlisting}
        If every operation requires a network round-trip, latency accumulates:
        \begin{equation*}
            T \approx n \cdot RTT
        \end{equation*}
        Where $T$ is the total time, $n$ is the number of items, and $RTT$ is the round-trip time (i.e., the time it takes for a message to go from $A$ to $B$ and back). So architecture and algorithms must take communication latency seriously.

    \item \important{Bandwidth is infinite}. \hl{A network can move only a finite amount of data per unit time:}
        \begin{equation*}
            \text{bandwidth} < \infty
        \end{equation*}
        If many machines exchange large amounts of information, the network itself may become the bottleneck.

        Imagine $1000$ machines constantly sending gigabytes of state to each other. Even if every CPU is mostly idle, the system can still perform badly because the \textbf{network capacity is a resource}. This is one reason distributed algorithms often try to reduce communication.

    \item \important{The network is secure}. \hl{Once communication crosses a network, additional attack surfaces appear.} Data may potentially be intercepted, modified, or an attacker may impersonate another component.

        So communication may require authentication, encryption, integrity protection, authorization. This fallacy connects directly with the security challenge that will appear later in the notes.

    \item \important{Topology doesn't change}. Suppose our system contains: $A$, $B$, $C$, and $D$ connected in a line. It is easy to design an algorithm assuming that this arrangement will remain stable. But \textbf{distributed systems change}. A node can \textbf{fail}:
        \begin{equation*}
            A-B-\cancel C-D
        \end{equation*}
        A new node can join:
        \begin{equation*}
            A-B-C-D-E
        \end{equation*}
        A mobile device can move, a router can disappear, or a service can migrate. Therefore the \hl{set of nodes and connections is dynamic.} This becomes especially obvious in the mobile and ad-hoc systems we saw in the previous section (\autopageref{def:ad-hoc}).

    \item \important{There is one administrator}. On our own computer, one administrative authority can generally configure the whole system. Instead, large distributed systems may cross organizations, departments, companies, cloud providers, countries. \hl{There may be no single administrator who controls everything.} Different domains may have different policies, security rules, configurations, software versions. So distributed systems often have to \textbf{operate across administrative boundaries}.

    \item \important{Transport cost is zero}. \hl{Communication has a cost.} That cost can include:
        \begin{equation*}
            \text{time} + \text{bandwidth} + \text{CPU} + \text{energy} + \text{possibly money}
        \end{equation*}
        Consider two implementations.
        \begin{itemize}
            \item \textbf{Version A}. Send $1000$ remote messages:
                \begin{equation*}
                    1000 \times \text{network operation}
                \end{equation*}
            \item \textbf{Version B}. Aggregate the information and send one message:
                \begin{equation*}
                    1 \times \text{network operation}
                \end{equation*}
        \end{itemize}
        Even if they compute logically the \textbf{same result}, their \textbf{performance} can be \textbf{dramatically different}. So \hl{communication is not free}. This idea will appear repeatedly throughout distributed systems.

    \item \important{The network is homogeneous}. \hl{A distributed system may combine machines with completely different characteristics}: x86 server, ARM phone, IoT sensor, cloud VM, using different hardware, operating systems, network technologies, data representations, programming languages. Therefore \hl{distributed systems are inherently heterogeneous}.
\end{enumerate}
\textcolor{Green3}{\faIcon{question-circle} \textbf{What the Eight Fallacies of Distributed Computing mean for us?}} We should not memorize them merely as 8 sentences. They are really \textbf{8 warnings} against pretending that:
\begin{equation*}
    \text{remote} \equiv \text{local}
\end{equation*}
A remote component differs fundamentally from a local one because the \textbf{network lies between us and it}. So a good distributed-system design starts from almost the opposite attitude: \textbf{assume that the network is unreliable, slow, and costly.} And then \hl{builds guarantees \textbf{on top of that unreliable and changing environment}.}

\highspace
\begin{takeawaysbox}[: Why Distributed Systems Are Special]
    A distributed system is fundamentally different from a centralized one
    because components communicate through a \textbf{network}:
    \begin{equation*}
        \boxed{
            \text{local interaction}
            \neq
            \text{remote interaction}
        }
    \end{equation*}
    A network introduces delay, limited capacity, failures, changing topology,
    security concerns, and heterogeneity. Consequently, distributed systems
    are:
    \begin{center}
        \textbf{more complex to design, harder to build, and much harder to
        build correctly.}
    \end{center}
    \important{The Eight Fallacies of Distributed Computing} summarize the
    assumptions that must \textbf{not} be made when designing a distributed
    system:
    \begin{enumerate}

        \item The network is reliable.

        \item Latency is zero.

        \item Bandwidth is infinite.

        \item The network is secure.

        \item Topology does not change.

        \item There is one administrator.

        \item Transport cost is zero.

        \item The network is homogeneous.

    \end{enumerate}
    They can be interpreted as four broad warnings:
    \begin{itemize}

        \item \important{Communication is uncertain}.
            Sending a message does not imply that it is received, and delays may
            be unpredictable.

        \item \important{Communication is not free}.
            Remote operations have latency, consume bandwidth, and require
            additional processing and resources.

        \item \important{The environment changes}.
            Nodes, links, locations, administrators, and network topology may
            change over time.

        \item \important{Distributed systems are heterogeneous}.
            Components may use different hardware, operating systems, networks,
            protocols, and software environments.

    \end{itemize}
    Therefore, a fundamental distributed-systems principle is:
    \begin{equation*}
        \boxed{
            \text{do not design a remote interaction as if it were local}
        }
    \end{equation*}
    Correct distributed software must explicitly deal with communication,
    concurrency, uncertainty, and failures instead of assuming them away.
\end{takeawaysbox}